\documentclass[11pt]{article}

\usepackage[margin=1in]{geometry}
\usepackage[T1]{fontenc}
\usepackage[utf8]{inputenc}
\usepackage{lmodern}
\usepackage{amsmath,amssymb,amsthm,mathtools}
\usepackage{booktabs,tabularx}
\usepackage{xcolor}
\usepackage{hyperref}
\usepackage{xurl}
\usepackage{microtype}

\hypersetup{colorlinks=true,linkcolor=blue,citecolor=blue,urlcolor=blue}
\setlength{\parindent}{0pt}
\setlength{\parskip}{0.55em}
\setlength{\emergencystretch}{3em}

\newcommand{\OPH}{\mathrm{OPH}}
\newcommand{\dd}{\,\mathrm d}
\newtheorem{theorem}{Theorem}[section]
\newtheorem{proposition}[theorem]{Proposition}
\newtheorem{corollary}[theorem]{Corollary}
\newtheorem{definition}[theorem]{Definition}

\title{Observer-Patch Holography and the Dark Matter Phenomenon\\
\large The Modular-Anomaly Source}
\author{Bernhard Mueller \and David Matscheko}
\date{\OPHPaperReleaseDate}
\newcommand{\OPHPaperReleaseID}{r2011}
\newcommand{\OPHPaperReleaseDate}{August 1, 2026}
\newcommand{\OPHPaperReleaseBanner}{%
\begin{center}
\small\textbf{Paper release:} \texttt{\OPHPaperReleaseID}\quad\textbf{Released:} \OPHPaperReleaseDate
\end{center}
}
\newcommand{\PaperReleaseID}{\OPHPaperReleaseID}
\newcommand{\PaperReleaseDate}{\OPHPaperReleaseDate}
\newcommand{\PaperReleaseBanner}{%
\begin{center}
\small\textbf{Paper release:} \texttt{\PaperReleaseID}\quad\textbf{Released:} \PaperReleaseDate
\end{center}
}


\begin{document}
\maketitle
\begin{center}
\small\textbf{Paper release:} \OPHPaperReleaseID\quad
\textbf{Released:} \OPHPaperReleaseDate
\end{center}

\begin{center}
\small\textbf{Author affiliation:} Bernhard Mueller, Pragma Research Inc.
\end{center}

\begin{abstract}
In observer-patch holography the Einstein equation sources curvature from the
modular-charge stress of the screen under a universal coupling. Every screen
source that carries modular charge therefore gravitates, whether or not it
carries the electromagnetic readout, and the luminous-only Einstein relation
holds exactly when the non-luminous charge vanishes. A dark sector is
generic. The candidate source is the anomalous modular energy carried on
overlap collars, the same finite-stage remainder that the Einstein derivation
controls. Its rest-frame stress is bounded by the collar recovery defect,
vanishes at full recovery, and dilutes as the inverse cube of the scale
factor when conserved. These statements are machine-checked. The source must
meet two data boundaries: on the SPARC rotation-curve sample the deep-regime
law reproduces the radial-acceleration and Tully--Fisher relations with one
constant near \(1.4\times10^{-10}\,\mathrm{m\,s^{-2}}\), and the Cassini
quadrupole excludes any universal fixed-coefficient extension of that law at
nineteen standard deviations. The magnitude of the anomalous charge, its
radial profile, and the vanishing of the recovery defect at high record
density are work in progress.
\end{abstract}

\section*{Claim Boundary}

The finite OPH results used here are the rest-frame Einstein relation on its
named small-ball premises, the decomposition of the cap modular generator into
a bulk part, a central part, and a collar-carried anomalous part, and the
controlled bound on that anomalous part~\cite{ophEinstein,ophmicro}. The
theorems of Sections~\ref{sec:generic} and \ref{sec:source} are algebraic
consequences of those objects and are checked in the Lean
formalization~\cite{ophlean}. Universal coupling is a premise of the Einstein
branch, named as a structure field in the formalization and not established
there.

What the paper does not supply: a quotient-invariant receipt localizing the
anomalous modular energy on collars, the value of the collar constant, the
radial profile of the anomalous charge around a baryonic source, a proof that
the recovery defect vanishes in any physical regime, and any cosmological
likelihood. No cosmological result carries an OPH falsification verdict.

\tableofcontents

\section{What Sources Gravity in OPH}

The Einstein arrow of the spacetime paper ends in
\begin{equation}
G_{ab}+\Lambda g_{ab}=8\pi G\langle T_{ab}\rangle,
\label{eq:einstein}
\end{equation}
where \(\langle T_{ab}\rangle\) is the local conserved stress reconstructed
from positive null translations and modular charges of the record
algebra~\cite{ophEinstein}. In the rest frame the relation is obtained from
three premises on a small geodesic ball of radius \(\ell\): generalized
entropy is stationary, the bulk entropy variation equals
\((8\pi^2\ell^4/15)\,\delta\langle T_{00}\rangle\), and the fixed-volume area
variation equals \(-(4\pi\ell^4/15)\,\delta G_{00}\). Exact arithmetic then
gives \(\delta G_{00}=8\pi G\,\delta\langle T_{00}\rangle\).

The stress entering this relation is modular charge. The screen does not
consult the electromagnetic readout of a source before responding to it. That
is the content of the universal-coupling premise, and it is the only input
the dark-sector question needs from the gravity derivation.

\section{The Generic Dark Sector}
\label{sec:generic}

Let a finite family of screen sources be indexed by \(i\), each with
rest-frame modular-charge contribution \(q_i\) and a flag recording whether
it carries the electromagnetic readout. Write \(Q_{\rm lum}\) for the sum over
flagged sources and \(Q_{\rm dark}\) for the sum over the others, so that
\(Q_{\rm tot}=Q_{\rm lum}+Q_{\rm dark}\).

\begin{definition}[Universal coupling]
The stress contraction entering the rest-frame Einstein relation equals
\(Q_{\rm tot}\), with no readout-dependent weight.
\end{definition}

\begin{theorem}[Geometry responds to all modular charge]
\label{thm:all-charge}
Under universal coupling,
\(\delta G_{00}=8\pi G\,(Q_{\rm lum}+Q_{\rm dark})\).
\end{theorem}

\begin{theorem}[The luminous-only relation]
\label{thm:iff}
Under universal coupling, \(\delta G_{00}=8\pi G\,Q_{\rm lum}\) holds if and
only if \(Q_{\rm dark}=0\).
\end{theorem}

\begin{corollary}[A geometric excess forces a dark source]
\label{cor:excess}
If \(\delta G_{00}\neq8\pi G\,Q_{\rm lum}\), some source without the
electromagnetic readout carries nonzero modular charge.
\end{corollary}

\begin{proposition}[Attractive sign]
If every non-luminous charge is nonnegative, \(Q_{\rm dark}\ge0\), and the
dark contribution enters the weak-field equation with the sign of baryonic
mass.
\end{proposition}

The proofs are finite sums and the positivity of \(8\pi G\). Their value is
the direction of the inference. In a theory where gravity is sourced by
modular charge, the question is not whether a dark sector exists; any
geometric excess over the luminous prediction is by itself the statement
that non-luminous modular charge is present. The open questions are which
screen object carries that charge, how much of it there is, and how it is
distributed.

\section{The Modular-Anomaly Source}
\label{sec:source}

On the support-visible branch the cap modular generator has the local form
\begin{equation}
K_C=2\pi B_C+K_C^{\rm anom}+\mathrm{const},
\label{eq:split}
\end{equation}
with \(B_C\) the bulk boost generator and \(K_C^{\rm anom}\) the nonadditive
part carried on the overlap collar~\cite{ophEinstein}. At finite regulator
stage the Einstein derivation carries this part as a remainder and bounds it:
\begin{equation}
\bigl|\delta\langle E_C^{(\delta)}\rangle\bigr|\le C_C\,\eta_\delta,
\label{eq:bound}
\end{equation}
where \(\eta_\delta\) is the collar recovery defect and \(C_C\) depends only
on the collar model and the bounded observable class.

\begin{definition}[Anomalous stress]
\begin{equation}
\delta\langle T_{00}^{\rm anom}\rangle
:=\frac{15}{8\pi^2\ell^4}\,\delta\langle E_C^{(\delta)}\rangle .
\label{eq:tanom}
\end{equation}
\end{definition}

This is the Einstein branch's own bookkeeping term. The dark-sector proposal
is that it is not bookkeeping: the anomalous modular energy is a physical
non-luminous charge in the sense of Section~\ref{sec:generic}. Three
consequences follow from \eqref{eq:bound} and \eqref{eq:tanom} alone.

\begin{theorem}[Bounded source]
\label{thm:bounded}
\(\bigl|\delta\langle T_{00}^{\rm anom}\rangle\bigr|
\le 15\,C_C\,\eta_\delta/(8\pi^2\ell^4)\).
\end{theorem}

\begin{theorem}[Full recovery switches the source off]
\label{thm:off}
If \(\eta_\delta=0\) then \(\delta\langle T_{00}^{\rm anom}\rangle=0\).
Along any family with fixed \(\ell\), bounded \(C_C\), and
\(\eta_\delta\to0\), the anomalous stress tends to zero.
\end{theorem}

\begin{theorem}[Cold scaling]
\label{thm:cold}
A conserved anomalous charge in a comoving cell of physical volume
\(\propto a^3\) has density \(\propto a^{-3}\).
\end{theorem}

Theorem~\ref{thm:off} is the screening interface. Wherever repair saturates
and the collar recovery defect vanishes, the dark source is absent. Whether
\(\eta_\delta\) vanishes in the high-acceleration, high-record-density regime
of a planetary system is the physical question on which the Solar-System
bound below turns; the theorem states what follows if it does.

\section{Data Boundaries the Source Must Meet}
\label{sec:data}

\subsection{Deep regime}

A static spherical source law of the form
\begin{equation}
a_{\rm dark}=\sqrt{a_b\,a_0},
\qquad
v^4=GM_b\,a_0,
\label{eq:deep}
\end{equation}
is the deep-regime target. One OPH mechanism for its exponent is a
codimension-one collar count: repair opportunities on closed cuts around a
point source scale as \(\sqrt{g_b/a_0}=r_M/r\) rather than as the area count
\(g_b/a_0\)~\cite{ophmicro}. That count is a premise, not a consequence of
Section~\ref{sec:source}.

The law is compared with the SPARC sample of 175 disk galaxies with
\(3.6\,\mu\)m photometry~\cite{sparc2016} under the standard cuts of McGaugh,
Lelli, and Schombert~\cite{mls2016}: quality flag at most 2, inclination at
least \(30^\circ\), velocity error at most ten percent, stellar
mass-to-light ratios \(0.5\) for disks and \(0.7\) for bulges, leaving 2700
points in 149 galaxies~\cite{ophrar}. On the self-consistent deep subset
\(g_{\rm bar}<fa_0\) the fixed-exponent fit gives
\[
a_0=1.57\times10^{-10},\quad
1.40\times10^{-10},\quad
1.28\times10^{-10}\ \mathrm{m\,s^{-2}}
\]
for \(f=0.3\), \(0.1\), \(0.03\), with 1651, 1043, 233 points and rms
scatter between \(0.14\) and \(0.17\) dex; the free exponents are
\(0.60\pm0.01\), \(0.57\pm0.02\), \(0.39\pm0.08\). For the 123 galaxies with
a measured flat velocity, \(M_b=0.5L_{[3.6]}+1.33M_{\rm HI}\), the
fixed-exponent Tully--Fisher normalization gives
\(a_0=1.55\times10^{-10}\,\mathrm{m\,s^{-2}}\) with \(0.26\) dex scatter and
free exponent \(3.75\pm0.10\). The two constants differ by \(0.046\) dex,
inside the combined scatter. The exponents are the parameter-free content of
\eqref{eq:deep}; \(a_0\) is a fitted comparison value, and no OPH
construction supplies it.

\subsection{Solar System}

Applied unchanged to the Sun inside the measured Galactic field, the
collar-count law with the fixed coefficient \(1-P/24\) predicts an
external-field quadrupole
\(Q_2=3.62\times10^{-26}\,\mathrm{s^{-2}}\), against the Cassini value
\((1.6\pm1.8)\times10^{-27}\,\mathrm{s^{-2}}\)~\cite{park2026}: a
\(19\sigma\) exclusion at fixed inputs, and about \(11\sigma\) after the Gaia
uncertainty on the Galactic acceleration~\cite{gaia2021}. Equivalently, the
additive form \(g=g_{\rm bar}+\sqrt{g_{\rm bar}a_0}\) predicts a relative
anomaly of \(1.5\times10^{-4}\) at one astronomical unit against a bound near
\(10^{-8}\). Any universal fixed-coefficient extension of \eqref{eq:deep} is
excluded.

This is why Theorem~\ref{thm:off} is the first open target rather than a
remark. A suppression rule introduced after reading Cassini would be a fitted
continuation. A proof that the collar recovery defect vanishes where record
density saturates would be the same object that sources the halo switching
itself off, and it would be decided before any further data are examined.

\section{Open Targets}
\label{sec:open}

\begin{enumerate}
\item \textbf{Screening from recovery.} Prove or refute that \(\eta_\delta\to0\)
in the repair-saturated regime of a compact source, with a rate in the local
acceleration or record density. Theorem~\ref{thm:off} then gives the
Solar-System suppression; a refutation retires the modular-anomaly source.
\item \textbf{Magnitude.} Attach the normalization of
\(\delta\langle E_C^{(\delta)}\rangle\) to repair conservation on the
Ward/Bianchi rung of the gravity premise ladder, so that \(a_0\) is a
consequence rather than a value borrowed from \(H_0\) and \(\Lambda\).
\item \textbf{Profile.} Derive the codimension-one collar count, or replace
it, from the collar geometry of a point source, so that the exponent in
\eqref{eq:deep} is a theorem on the same premises as
Section~\ref{sec:source}.
\item \textbf{Localization.} Supply the quotient-invariant receipt identifying
the collar operator in \eqref{eq:split} with the small-ball remainder in
\eqref{eq:tanom}.
\end{enumerate}

A rotor completion, with integer repair occupation and a conjugate phase,
is admissible only under the identification of the occupation with the
anomalous modular charge and of the phase with the modular flow angle. Without
that identification it is a parametrization of items 2 and 3 rather than a
derivation, and it is not used here.

\section{Conclusion}

Under universal coupling the screen gravitates according to its total
modular charge. A dark sector is therefore a theorem-level feature of OPH,
and any geometric excess over the luminous prediction is the presence of
non-luminous modular charge. The anomalous modular energy on overlap collars
is the one candidate for that charge that the Einstein derivation itself
carries; it is bounded by the recovery defect, absent at full recovery, and
cold when conserved. The deep-regime galaxy data fix the target its profile
must reach, and the Cassini bound fixes the target its screening must reach.
The proofs that close those targets are work in progress.

\begin{thebibliography}{99}

\bibitem{ophEinstein}
B. Mueller et al.,
\emph{Recovering Observer Spacetime and Einstein Dynamics from Overlap Consistency},
OPH spacetime and Einstein paper.
\url{https://github.com/FloatingPragma/observer-patch-holography/blob/main/paper/recovering_observer_spacetime_and_einstein_dynamics_from_overlap_consistency.pdf}

\bibitem{ophmicro}
B. Mueller et al.,
\emph{Federated Echosahedral Screen Microphysics: Patch Hardware, Records, and Observer Synchronization in OPH}.
\url{https://github.com/FloatingPragma/observer-patch-holography/blob/main/paper/screen_microphysics_and_observer_synchronization.pdf}

\bibitem{ophlean}
B. Mueller et al.,
\emph{Observer-Patch Holography: machine-checked Lean formalization},
Einstein-branch modules \texttt{SmallBall} and \texttt{DarkSector},
2026.
\url{https://github.com/FloatingPragma/observer-patch-holography/tree/main/Lean/ObserverPatchHolography/EinsteinBranch}

\bibitem{ophrar}
B. Mueller,
\emph{Deep-regime comparison of the OPH galaxy law with SPARC},
2026.
\url{https://github.com/FloatingPragma/observer-patch-holography/tree/main/code/cosmology/rar_deep_regime}

\bibitem{sparc2016}
F. Lelli, S. S. McGaugh, and J. M. Schombert,
\emph{SPARC: Mass Models for 175 Disk Galaxies with Spitzer Photometry and Accurate Rotation Curves},
Astron. J. 152, 157 (2016), arXiv:1606.09251.
\url{https://arxiv.org/abs/1606.09251}

\bibitem{mls2016}
S. S. McGaugh, F. Lelli, and J. M. Schombert,
\emph{Radial Acceleration Relation in Rotationally Supported Galaxies},
Phys. Rev. Lett. 117, 201101 (2016), arXiv:1609.05917.
\url{https://arxiv.org/abs/1609.05917}

\bibitem{park2026}
R. S. Park et al.,
\emph{Cassini spacecraft navigation data constrain modified gravity in the Solar System},
2026, arXiv:2602.17884.
\url{https://arxiv.org/abs/2602.17884}

\bibitem{gaia2021}
Gaia Collaboration,
\emph{Gaia Early Data Release 3: Acceleration of the Solar System from Gaia astrometry},
Astron. Astrophys. 649, A9 (2021), arXiv:2012.02036.
\url{https://arxiv.org/abs/2012.02036}

\end{thebibliography}

\section*{AI Assistance Disclosure}
This research project used research-grade commercial models, including
Anthropic's Fable and OpenAI's GPT-5.6-Sol, for research support, software
development, editing, and synthesis. The authors are responsible for the
paper's claims, methods, and final text.

\end{document}
